\section{Technical Validation}
To validate the ``AI-readiness'' of the dataset, we trained baseline deep learning models for two complementary tasks: 2D lung ROI extraction and 3D nodule segmentation. Together, these form a two-stage inference pipeline, in which the 2D model provides a lung bounding box and the 3D segmentation is performed within this region.
All reported metrics were computed exclusively on the held-out test split described earlier, which was never used for any other purpose.

The 1,980-series training pool was further partitioned into a training set (1,683 series) and a validation set (297 series), using the same methodology we described in connection with the test split.
The validation series were used only for checkpoint selection: from the periodic checkpoints of each run we retained the one with the highest mean per-series (for the 2D task: per-slice) Dice coefficient on the validation set.
We note that the released dataset fixes only the test split.
However, the validation assignment we used is distributed with the training code so the baselines can be reproduced exactly.

\subsection{Notation and evaluation metrics}
\label{sec:metrics}
Let $\mathcal{S}$ denote the set of test series. Each series $s \in \mathcal{S}$ is defined on its own voxel grid $\Omega_s \subset \mathbb{Z}^3$, whose extent varies from series to series (the shared quantity is the isotropic 1\,mm voxel size, not the grid shape) (for the 2D task, each axial slice is defined on a pixel grid $\Omega_s \subset \mathbb{Z}^2$, and ``series'' below should be read as ``slice''). We write $Y_s \subseteq \Omega_s$ for the set of ground-truth foreground voxels of $s$ (nodule voxels in the 3D task, lung pixels in the 2D task) and $\hat{Y}_s \subseteq \Omega_s$ for the set of voxels predicted as foreground.

Lung-nodule segmentation is an extremely imbalanced problem: even within the lung-cropped volumes used for the 3D task, nodule voxels account for fewer than one voxel in a thousand, so voxel accuracy is unusable as a metric. We therefore report three complementary sets of metrics.

\paragraph{Micro-averaged voxel metrics.}
Confusion counts are pooled over the whole test split,
\begin{equation*}
\mathrm{TP} = \sum_{s \in \mathcal{S}} |Y_s \cap \hat{Y}_s|, \qquad
\mathrm{FP} = \sum_{s \in \mathcal{S}} |\hat{Y}_s \setminus Y_s|, \qquad
\mathrm{FN} = \sum_{s \in \mathcal{S}} |Y_s \setminus \hat{Y}_s|,
\end{equation*}
and precision, recall, and the Dice coefficient (equivalently, the F1-score) are computed from the pooled counts:
\begin{equation*}
\mathrm{Prec} = \frac{\mathrm{TP}}{\mathrm{TP}+\mathrm{FP}}, \qquad
\mathrm{Rec} = \frac{\mathrm{TP}}{\mathrm{TP}+\mathrm{FN}}, \qquad
\mathrm{Dice} = \frac{2\,\mathrm{TP}}{2\,\mathrm{TP}+\mathrm{FP}+\mathrm{FN}}.
\end{equation*}
Micro-averaging weights every voxel equally, so large lesions dominate the score.

\paragraph{Per-case Dice.}
To weight every patient equally regardless of lesion size, we also report the mean of the per-series Dice coefficients
\begin{equation*}
\mathrm{Dice}_s = \frac{2\,|Y_s \cap \hat{Y}_s|}{|Y_s| + |\hat{Y}_s|},
\end{equation*}
with the convention $\mathrm{Dice}_s = 1$ when $Y_s = \hat{Y}_s = \emptyset$.

\paragraph{Instance-level metrics.}
Let $\mathcal{C}(Y_s)$ denote the set of 26-connected components of $Y_s$, each component representing one nodule instance. Instance recall is the fraction of ground-truth components touched by at least one predicted voxel, and instance precision is the fraction of predicted components touching at least one ground-truth voxel, both pooled over the split:
\begin{equation*}
\mathrm{IR} = \frac{\sum_{s} \bigl|\{\, c \in \mathcal{C}(Y_s) : c \cap \hat{Y}_s \neq \emptyset \,\}\bigr|}{\sum_{s} |\mathcal{C}(Y_s)|},
\qquad
\mathrm{IP} = \frac{\sum_{s} \bigl|\{\, \hat{c} \in \mathcal{C}(\hat{Y}_s) : \hat{c} \cap Y_s \neq \emptyset \,\}\bigr|}{\sum_{s} |\mathcal{C}(\hat{Y}_s)|}.
\end{equation*}
Note that a single overlapping voxel already counts as a hit and no one-to-one matching between predicted and ground-truth components is enforced; these are deliberately permissive detection criteria rather than IoU-thresholded instance-segmentation scores. No minimum-size filtering is applied to predicted components: every isolated false-positive component counts, including single-voxel ones. We deliberately report these raw values; the choice of a component-size threshold or of a candidate-classification stage is a downstream design decision that we leave to users of the dataset.

\paragraph{Patch-based training and sliding-window inference.}
All baselines are trained on fixed-size patches at the native 1\,mm resolution and applied with sliding-window inference; resizing a lung crop or an axial slice to a fixed grid would distort the morphology of the lung and of small nodules.
During training, each training volume contributes four randomly placed patches per epoch; with probability $2/3$ the patch is centred on a uniformly chosen nodule voxel and with probability $1/3$ on a voxel chosen uniformly from the non-air region of the volume (voxels with normalised intensity above zero), which biases the patch distribution towards nodule-containing patches without excluding pure-background ones (for the 2D task the two probabilities are $1/2$ and $1/2$, with lung pixels as the positive class).
At inference, windows of the training patch size are placed on a regular grid with a stride of half the window size along \emph{each} axis. The overlap fraction is therefore a per-axis notion: two windows adjacent along a single axis share $50\%$ of their volume, windows offset along two or all three axes share $25\%$ and $12.5\%$, respectively, and every voxel in the interior of the volume is covered by $2^3 = 8$ windows ($2^2 = 4$ windows in the 2D task).
Each window $w$ predicts a full vector of class logits at every voxel it covers; we write $z_{w,c}(v)$ for the logit of class $c$ predicted by window $w$ at voxel $v$. The per-window logits are blended into a single full-resolution logit field, one class channel at a time:
\begin{equation*}
z_c(v) = \frac{\sum_{w \ni v} g_w(v)\, z_{w,c}(v)}{\sum_{w \ni v} g_w(v)},
\end{equation*}
where the sum runs over all windows $w$ containing $v$ and $g_w$ is an isotropic Gaussian importance map centred on the window centre with standard deviation $0.125$ times the window size. The Gaussian weighting down-weights window-border predictions, which see truncated context, in favour of window-centre predictions, and removes seam artefacts at window boundaries. The final mask assigns to voxel $v$ the class $\arg\max_c z_c(v)$ — the arg-max is taken over the class index at each voxel separately, never over voxels — which in the binary case is equivalent to a $0.5$ threshold on the softmax foreground probability (for the single-channel 2D task, a $0.5$ threshold on the blended sigmoid output). The number of foreground voxels is therefore unconstrained.

All models were trained on a single NVIDIA H100 GPU; the complete configuration of every run (optimiser, schedule, augmentation, seed) is distributed with the training code.
Note that these baselines serve to verify the dataset rather than to optimise segmentation performance.

\subsection{2D ROI Segmentation}
Tightly cropping the CT volume to the lung region before 3D nodule segmentation reduces the computational cost of the 3D stage and removes anatomy that cannot contain a nodule. We trained two 2D baseline models for this purpose, SegResNet \cite{Myronenko2019} (6.9M parameters) and a small SwinUNETR \cite{Hatamizadeh2022} variant (feature size 24, 6.3M parameters, chosen to match the parameter budget of the SegResNet).
Both operate on individual axial slices at the native 1~mm resolution and produce a binary lung mask.
Training uses $256 \times 256$ patches sampled as described in Section~\ref{sec:metrics}; inference uses $256 \times 256$ sliding windows with the $50\%$-per-axis overlap and Gaussian blending defined there.
At inference time, the union of per-slice lung masks defines the 3D lung ROI.
The models were trained with a soft Dice loss on the sigmoid output (with squared class cardinalities in the denominator \cite{Milletari2016}) on the 2D slices of the training-split series that carry a lung annotation (365,014 slices from NSCLC-Radiomics and NLST) and evaluated on the 45,751 slices of the corresponding test-split series.

Table~\ref{tab:roi2d} summarises the results. Both models are close to the ceiling of this task, whose target is a large, contrast-rich region of simple morphology, and the two architectures are within a tenth of a Dice point of each other on every micro-averaged metric. SwinUNETR is nevertheless ahead on every reported metric, and the difference concentrates in the tail of the per-slice Dice distribution: at the 5th percentile SwinUNETR retains a Dice of $0.8692$ against the SegResNet's $0.8352$. The low-scoring slices are those where the lung is absent or barely present (the top and bottom of the stack), where a few false-positive pixels can drive the per-slice Dice towards zero; the global receptive field of the transformer handles these anatomically ambiguous slices somewhat more conservatively. We therefore use the SwinUNETR model to compute the ROI for the 3D stage.

\begin{table}
\caption{2D lung ROI segmentation on the held-out test split (45,751 axial slices). Micro-averaged metrics and per-slice statistics as defined in Section~\ref{sec:metrics}; the per-slice statistics weight every slice equally.}
\label{tab:roi2d}
\centering
\small
\setlength{\tabcolsep}{4pt}
\begin{tabular}{lccccccc}
\toprule
& & \multicolumn{3}{c}{Micro-averaged} & \multicolumn{2}{c}{Per-slice Dice} \\
\cmidrule(lr){3-5} \cmidrule(lr){6-7}
Model & Parameters & Dice & Precision & Recall & Mean & p05 \\
\midrule
SegResNet & 6.9M & 0.9823 & 0.9832 & 0.9814 & 0.9511 & 0.8352 \\
SwinUNETR & 6.3M & \textbf{0.9834} & \textbf{0.9836} & \textbf{0.9832} & \textbf{0.9612} & \textbf{0.8692} \\
\bottomrule
\end{tabular}
\end{table}

\subsection{3D Nodule Segmentation}
We trained three 3D baseline models on lung-cropped volumes: a small SegResNet (initial feature width of 16, 20.7M parameters), a wide SegResNet (initial feature width of 32, 82.7M parameters), and a DynUNet, an nnU-Net-style 3D U-Net implemented in MONAI \cite{MONAIDynUNet} (16.5M parameters).

Each series was cropped to its lung bounding box, padded by 20 voxels in every direction, and kept at the native 1~mm voxel grid; no resizing was applied. For training, the bounding boxes were derived from the ground-truth lung masks for NSCLC-Radiomics and NLST and, for LIDC-IDRI, which has no lung annotation, from the per-slice lung predictions of a 2D SegResNet lung model trained on the NSCLC-Radiomics and NLST lung annotations (the predecessor of the SegResNet baseline of the previous subsection).
At inference on a new scan, the bounding box comes from the 2D ROI model described above.

The models were trained on $128 \times 128 \times 128$ patches with the nodule-biased sampling of Section~\ref{sec:metrics}, and applied with $128^3$ sliding windows, $50\%$-per-axis overlap, and Gaussian blending as defined there.
The training loss combines a focal Tversky term \cite{Abraham2019} with a lightly weighted cross-entropy term. Writing $p_v \in [0,1]$ for the predicted nodule probability at voxel $v$ (the softmax output) and $y_v \in \{0,1\}$ for the label, the soft confusion counts of a patch are
\begin{equation*}
\mathrm{tp} = \sum_v p_v y_v, \qquad
\mathrm{fp} = \sum_v p_v (1 - y_v), \qquad
\mathrm{fn} = \sum_v (1 - p_v)\, y_v,
\end{equation*}
and the loss is
\begin{equation*}
\mathcal{L} = \bigl(1 - \mathrm{TI}_{\alpha,\beta}\bigr)^{\gamma} + \lambda\, \mathcal{L}_{\mathrm{CE}},
\qquad
\mathrm{TI}_{\alpha,\beta} = \frac{\mathrm{tp} + \varepsilon}{\mathrm{tp} + \alpha\,\mathrm{fp} + \beta\,\mathrm{fn} + \varepsilon},
\end{equation*}
with $\alpha = 0.3$, $\beta = 0.7$ (penalising false negatives more than false positives), focal exponent $\gamma = 2$, and $\varepsilon = 10^{-6}$; in the Tversky term the probabilities are additionally clamped below at $10^{-4}$ for numerical stability. The cross-entropy term is the class-weighted cross-entropy
\begin{equation*}
\mathcal{L}_{\mathrm{CE}} = -\,\frac{\sum_v w_{y_v} \log q_v}{\sum_v w_{y_v}},
\qquad
q_v = y_v\, p_v + (1 - y_v)(1 - p_v),
\end{equation*}
where $q_v$ is the predicted probability of the true class at voxel $v$ and the class weights are $(w_0, w_1) = (1, 10)$ for (background, nodule); the normalisation by the sum of the applied weights is the standard weighted-mean convention. The mixing weight is $\lambda = 0.1$ for the small SegResNet and the DynUNet and $\lambda = 0.3$ for the wide SegResNet; the small cross-entropy term provides a non-vanishing per-voxel gradient when the network predicts all-background, the classic failure mode of pure overlap losses.
The best checkpoint of each run was selected by the mean per-series validation Dice.

Table~\ref{tab:nodule3d} summarises the results on the 325 held-out test series. The wide SegResNet is the strongest model on every voxel-level metric and on per-case Dice, and the DynUNet, the smallest of the three, trails on precision in particular. The gap between the micro-averaged and the per-case Dice of every model shows that large nodules dominate the voxel-weighted score; the per-case distributions are wide (standard deviation 0.26--0.27 for all three models, with a per-case median of 0.46 for the wide SegResNet), reflecting the difficulty of the many small nodules in the corpus. At the instance level, all three models detect 86--87\% of the ground-truth nodule components, whereas instance precision is low (0.09--0.12) because, in the absence of any minimum-size filtering, every isolated false-positive component counts, including single-voxel ones (see Section~\ref{sec:metrics}).

\begin{table}
\caption{3D nodule segmentation on the held-out test split (325 series, 947 ground-truth nodule components). Metrics as defined in Section~\ref{sec:metrics}.}
\label{tab:nodule3d}
\centering
\small
\setlength{\tabcolsep}{4pt}
\begin{tabular}{lccccccc}
\toprule
& & \multicolumn{3}{c}{Micro-averaged} & Per-case & \multicolumn{2}{c}{Instance level} \\
\cmidrule(lr){3-5} \cmidrule(lr){6-6} \cmidrule(lr){7-8}
Model & Parameters & Dice & Precision & Recall & Dice (mean) & Recall & Precision \\
\midrule
SegResNet (small) & 20.7M & 0.5370 & 0.4645 & 0.6363 & 0.4261 & \textbf{0.8680} & 0.0863 \\
SegResNet (wide)  & 82.7M & \textbf{0.6064} & \textbf{0.4926} & \textbf{0.7886} & \textbf{0.4606} & 0.8648 & \textbf{0.1225} \\
DynUNet (3D U-Net) & 16.5M & 0.4966 & 0.3634 & 0.7839 & 0.4158 & 0.8659 & 0.1092 \\
\bottomrule
\end{tabular}
\end{table}
